% Chapter 3

\chapter{Desenho da Solução} % Main chapter title
\label{chap:Chapter3} % For referencing the chapter elsewhere, use \ref{chap:Chapter3}

%----------------------------------------------------------------------------------------
\section{Requisitos}
\label{sec:chap3_requirements}

Os requisitos resultam do levantamento feito com a Adhara. % TODO: confirmar/atualizar com a Adhara.

\subsection{Requisitos Funcionais}

\begin{itemize}
    \item \textbf{RF1.} O assistente responde a dúvidas sobre a pergunta em causa (enunciado, opções, explicações, dicas e informação-chave).
    \item \textbf{RF2.} O assistente usa exclusivamente a informação guardada na base de dados para essa pergunta, sem acesso à Internet nem a conhecimento externo.
    \item \textbf{RF3.} Se a informação necessária não estiver no contexto, o assistente di-lo claramente.
    \item \textbf{RF4.} Perguntas fora do âmbito da pergunta são recusadas com uma frase fixa.
    \item \textbf{RF5.} O assistente está disponível nos modos Teste e Estudo e não no modo Simulado.
    \item \textbf{RF6.} Quando a resposta correta não foi revelada ao aluno, o assistente não a revela.
    \item \textbf{RF7.} O assistente deve responder a perguntas escritas em várias línguas. % TODO: decidir a língua da resposta (sempre português ou língua do utilizador).
\end{itemize}

\subsection{Requisitos Não Funcionais}

\begin{itemize}
    \item \textbf{RNF1. Latência:} tempo de resposta aceitável entre 2 e 3 segundos.
    \item \textbf{RNF2. Custo:} custo mensal de modelos limitado a um orçamento definido com a Adhara.
    \item \textbf{RNF3. Escalabilidade:} 10 a 20\% dos utilizadores em simultâneo (cerca de 250 a 500 em 2500 utilizadores). % TODO: confirmar valores.
    \item \textbf{RNF4. Privacidade:} o envio dos dados das perguntas para uma \gls{API} externa foi indicado como aceitável pela Adhara; o tratamento de dados dos utilizadores deve cumprir o \gls{RGPD}.
    \item \textbf{RNF5. Validação:} respostas do assistente validadas pela equipa científica antes de serem consideradas fiáveis (processo a definir).
\end{itemize}

%----------------------------------------------------------------------------------------
\section{Contexto de Cada Pergunta}
\label{sec:chap3_context}

Cada pergunta do QBank está guardada com campos estruturados: vinheta clínica, opções de resposta, informação-chave, dicas, resposta correta, explicação (geral e de cada opção), catalogação e informação complementar. Para cada pedido, o contexto fornecido ao \gls{LLM} é construído a partir destes campos, em função do modo do QBank e de a resposta já ter sido revelada ou não.

% TODO: tabela com os campos incluídos no contexto em cada situação (Teste/Estudo com resposta revelada; sem resposta revelada).

%----------------------------------------------------------------------------------------
\section{Arquitetura}
\label{sec:chap3_architecture}

A Figura~\ref{fig:architecture} apresenta a arquitetura proposta. Antes de chegar ao \gls{LLM}, a pergunta do utilizador passa por um filtro de âmbito; depois da geração, a resposta é verificada antes de ser mostrada. % TODO: ajustar à arquitetura final.

\begin{figure}[t]
\centering
\begin{tikzpicture}[
  node distance=0.5cm,
  box/.style={draw, rounded corners, align=center, minimum height=1.1cm, text width=2.3cm, font=\small},
  arr/.style={-{Stealth[length=2mm]}}
]
\node[box] (user) {Pergunta do utilizador};
\node[box, right=of user] (filter) {Filtro de âmbito};
\node[box, right=of filter] (prompt) {Construção do \emph{prompt}};
\node[box, right=of prompt] (llm) {\gls{LLM}};
\node[box, below=0.9cm of llm] (check) {Verificação da resposta};
\node[box, left=of check] (answer) {Resposta ao utilizador};
\node[box, above=0.9cm of prompt, fill=gray!15] (db) {Dados da pergunta (BD)};

\draw[arr] (user) -- (filter);
\draw[arr] (filter) -- (prompt);
\draw[arr] (prompt) -- (llm);
\draw[arr] (llm) -- (check);
\draw[arr] (check) -- (answer);
\draw[arr] (db) -- (prompt);
\end{tikzpicture}
\caption[Arquitetura do assistente]{Arquitetura proposta do assistente conversacional.}
\label{fig:architecture}
\end{figure}

%----------------------------------------------------------------------------------------
\section{Estratégias de \emph{Prompting} e Restrições}
\label{sec:chap3_prompting}

O \emph{prompt} do assistente define o seu papel, as regras de contexto, de âmbito, de segurança e de estilo, e o formato de saída. As regras principais são:
\begin{itemize}
    \item usar apenas a informação do contexto e dizer quando esta não existe;
    \item recusar assuntos fora da pergunta com uma frase fixa;
    \item ignorar instruções do utilizador que tentem alterar as regras ou revelar o \emph{prompt};
    \item não escrever as letras das opções, referindo-se às opções pelo seu texto, porque a ordem das opções muda entre sessões.
\end{itemize}

% TODO: descrever as alterações ao prompt feitas a partir da avaliação (por exemplo, exemplos no prompt, regras repetidas no fim, saída estruturada com campo de âmbito).

%----------------------------------------------------------------------------------------
\section{Verificações Fora do Modelo}
\label{sec:chap3_guards}

Como as restrições só no \emph{prompt} podem não ser cumpridas por todos os modelos, a arquitetura inclui verificações determinísticas:
\begin{itemize}
    \item \textbf{antes do modelo:} filtro de âmbito para recusar perguntas claramente fora do tema;
    \item \textbf{depois do modelo:} deteção de fuga de instruções, de letras de opções e de valores que não existem no contexto.
\end{itemize}
% TODO: descrever a implementação e justificar com a avaliação (Capítulo 4).

%----------------------------------------------------------------------------------------
\section{Escolha do Modelo}
\label{sec:chap3_model}

A escolha entre um modelo executado localmente e um serviço externo depende da latência, do custo, da privacidade e da qualidade (fidelidade e cumprimento das regras). Os modelos são comparados no Capítulo~\ref{chap:Chapter4}.
% TODO: justificar a escolha final.